%{
%Formattting
\documentclass[11pt,a4paper]{article}
\usepackage[right=0.7in, left=0.7in, top=1in, bottom=0.75in]{geometry}
\usepackage{fancyheadings}
\usepackage{setspace}
\usepackage[T1]{fontenc}
\usepackage[utf8]{inputenc}
\setlength{\parindent}{0in}
\usepackage{listings}
%\usepackage{autobreak}
\lstset
{ %Formatting for code in appendix
	basicstyle=\footnotesize,
	numbers=left,
	stepnumber=1,
	showstringspaces=false,
	tabsize=1,
	breaklines=true,
	breakatwhitespace=false
}
\pagenumbering{gobble}
%Mathematics
\usepackage{amsmath}
\usepackage{amssymb}
\usepackage{mathtools}
\usepackage{float}
\usepackage{aliascnt}

\makeatletter
\newcommand{\crossout}[1]{%
	\begingroup
	\settowidth{\dimen@}{#1}%
	\setlength{\unitlength}{0.05\dimen@}%
	\settoheight{\dimen@}{#1}%
	\count@=\dimen@
	\divide\count@ by \unitlength
	\begin{picture}(0,0)
		\put(0,0){\line(20,\count@){20}}
		\put(0,\count@){\line(20,-\count@){20}}
	\end{picture}%
	#1%
	\endgroup
}
\usepackage{tikz}
\newcommand\halmos{\rule{.36em}{2ex}}
\newcommand{\powerset}[1]{\mathcal{P}(#1)}
\def\contradict{\tikz[baseline, x=0.22em, y=0.22em, line width=0.032em]\draw (0,2.83)--(2.83,0) (0.71,3.54)--(3.54,0.71) (0,0.71)--(2.83,3.54) (0.71,0)--(3.54,2.83);}

%Citations
%\usepackage{cite}
\usepackage[backend=bibtex,style=verbose-trad2]{biblatex} %works really really well, but no MLA format
%\usepackage[backend=biber,style=mla]{biblatex} %Doesn't print all sources for some reason
%\addbibresource{Math-IA.bib}
%\bibliography{MathIA} 
%Symbols and Other
\newcommand{\figref}[1]{Fig. \ref{#1}}
\newcommand{\source}[1]{\caption*{Source: {#1}} }
\usepackage{hyperref}
\hypersetup{
	colorlinks,
	allcolors=black
	%linkcolor=black,
	%urlcolor=black
}
%\addtolength{\oddsidemargin}{-.875in}
%\addtolength{\evensidemargin}{-.875in}
%\addtolength{\textwidth}{1.75in}
%\addtolength{\topmargin}{-.875in}
%\addtolength{\textheight}{1.75in}
\onehalfspacing
%\doublespacing
\pagestyle{fancy}
%\theoremstyle{definition}
%\newtheorem{defn}{Definición}
%\newtheorem{reg}{Regla}
%\newtheorem{ejer}{EJERCICIO}
%\pagestyle{fancyplain}
\headheight 32pt
%\lhead{\yourname\ \vspace{0.1cm} \\ \course}
%\chead{\textbf{\Large PUMP II}}
%}
\usepackage{pdfpages}
\lhead{\yourname\ \vspace{0.1cm}}
\chead{\textbf{\subject}\\\psetnum}
\rhead {\href{mailto:adityakrao.work@gmail.com}{adityakrao.work@gmail.com} \\
\href{mailto:adi.rao@mail.utoronto.ca}{adi.rao@mail.utoronto.ca}\\ \href{tel:14168025648}{+1 (416)-802-5648}}
\newcommand\course{University of Toronto}
\newcommand\psetnum{Reflection: Quantum Interference and Entanglement} 
\newcommand\yourname{Aditya Kanuru Rao}  
\newcommand{\subject}{\large{\course}}
\begin{document}
	\title{\vspace{-4cm}\psetnum}
	 %\title{\psetnum}  
        \author{\yourname}
	\date{\vspace{-0.3em}\textbf{Email:} \href{mailto:adityakrao.work@gmail.com}{adityakrao.work@gmail.com} \textbf{|} \href{mailto:adi.rao@mail.utoronto.ca}{adi.rao@mail.utoronto.ca}\\\textbf{Phone:} \href{tel:14168025648}{+1 (416)-802-5648}}
	%\titlepic{\includegraphics[width=0.6\linewidth]{images/lm-logo.png}}
        \vspace{0cm}
	\maketitle
\vspace{-1.2cm}
%\newpage
\rule{\linewidth}{0.5pt}

\textbf{}

The Quantum Interference and Entanglement experiment was one of the most interesting to me. It was a topic that I had a deep interest in from an interepretations of quantum mechanics point of view. Working with a lab partner was an adjustment, however, I felt that it did improve our overall efficiency in terms of getting results. However, it did restrict the amount of time I could be in the lab. Reason being that I did not want to inadvertaintly mess up the experimental setup while trying new things.

We did obtain relatively good results and I feel this ended up being one of hte most complete experiments I had done. Complete in the sense that my laboratory book was neat, well organized, time and date stamped, included uncertainty discussions for the final results, and included all procedures. Additionally, I included a lengthy disucssion about the interepretations of quantum mechanics and the possibility of extending this expreiment to replicate a `timing-experiment' (as done by Aspect).

That being said, there were some issues. The first being, in the begining of the experiment, I was not as careful with errorbars and uncertainty calculations. My justification was that we were just getting an initial feel for the setup. When asked about the errorbars of the `coincidence-values' during my experiment interview, I was only able to justify these as the standard deviation of the results. I do stand by this, however, I could have been better. Additionally, we had made a mistake with respect to the quartz-plate calibration. This may have adjusted our results, however, the overall experimental goal of showing entanglement was still achieved as we found $S=2.3\pm 0.5$ (uncertainty may have been smaller if we performed more runs).

I do have some minor issues with regards to the organization. However, I feel very satisfied with the work I completed within this experiment and proud of the lab notebook I produced.


\end{document}